\documentclass[12pt]{article}

\usepackage[spanish]{babel}
\usepackage[utf8]{inputenc}
\usepackage[T1]{fontenc}
\usepackage{geometry}
\usepackage{enumitem}
\usepackage{hyperref}
\usepackage{titlesec}
\usepackage{setspace}

\geometry{margin=2.5cm}
\setstretch{1.2}

\titleformat{\section}{\large\bfseries}{\thesection}{1em}{}

\title{\textbf{Guía de Práctica}\\
Manejo de Strings y Loops en LabVIEW}
\author{}
\date{}

\begin{document}

\maketitle
\thispagestyle{empty}

\section*{Objetivo General}

Comprender el manejo de \textit{strings}, estructuras iterativas y construcción dinámica de arreglos en LabVIEW, aplicándolo a un ejemplo progresivo que simula un flujo típico de adquisición y procesamiento de datos.

\newpage

\section{Programa 1 – Introducción a Strings}

\subsection*{Objetivo}
Familiarizarse con constantes, controles, indicadores y la función \textit{String Subset}.

\subsection*{Procedimiento}

\begin{enumerate}
    \item Crear una constante tipo \textbf{String}.
    \item Escribir el texto:
    \begin{center}
        \texttt{hello world}
    \end{center}
    \item Convertir la constante en control:
    \begin{itemize}
        \item Click derecho $\rightarrow$ \textit{Change to Control}
    \end{itemize}
    \item Insertar el bloque \textbf{String Subset}.
    \item Configurar:
    \begin{itemize}
        \item \texttt{offset = 0}
        \item \texttt{length = 1}
    \end{itemize}
    \item Conectar la salida a un indicador string.
\end{enumerate}

\subsection*{Resultado}

Se muestra únicamente la primera letra:

\begin{center}
\texttt{h}
\end{center}

Al modificar el \texttt{offset}, se pueden visualizar las demás letras.

---

\section{Programa 2 – Iteración con FOR Loop}

\subsection*{Objetivo}
Introducir estructuras iterativas y automatización según la longitud del string.

\subsection*{Procedimiento}

\begin{enumerate}
    \item Encerrar el \textit{String Subset} dentro de un \textbf{FOR Loop}.
    \item Conectar el iterador \texttt{i} al terminal \texttt{offset}.
    \item Mantener \texttt{length = 1}.
\end{enumerate}

\subsection*{Automatización}

\begin{enumerate}
    \item Insertar la función \textbf{String Length}.
    \item Conectarla al terminal \texttt{N} del FOR Loop.
\end{enumerate}

Ahora el programa funciona automáticamente para cualquier palabra.

---

\section{Programa 3 – Construcción de Array con Shift Register}

\subsection*{Objetivo}
Aprender a construir arreglos dinámicamente dentro de un loop.

\subsection*{Procedimiento}

\begin{enumerate}
    \item Insertar el bloque \textbf{Insert Into Array}.
    \item Agregar un \textbf{Shift Register} al FOR Loop.
    \item Conectar:
    \begin{itemize}
        \item Salida del array $\rightarrow$ lado derecho del shift register.
        \item Lado izquierdo $\rightarrow$ entrada superior del Insert Into Array.
    \end{itemize}
    \item Inicializar el shift register izquierdo con un array vacío:
    \begin{itemize}
        \item Click derecho $\rightarrow$ \textit{Create Constant}
    \end{itemize}
\end{enumerate}

\subsection*{Resultado}

Se obtiene un arreglo:

\begin{center}
\texttt{["h","e","l","l","o"," ","w","o","r","l","d"]}
\end{center}

El arreglo crece en tiempo real durante cada iteración.

---

\section{Programa 4 – Construcción con Auto-Indexing}

\subsection*{Objetivo}
Comparar métodos de construcción de arreglos.

\subsection*{Procedimiento}

\begin{enumerate}
    \item Conectar la letra directamente al borde del loop.
    \item Verificar que el túnel esté en modo \textbf{Auto-Indexing}.
\end{enumerate}

\subsection*{Comparación}

\begin{itemize}
    \item \textbf{Shift Register}: el arreglo crece durante la ejecución.
    \item \textbf{Auto-Index}: el arreglo se genera completo al finalizar el loop.
\end{itemize}

\subsection*{Aplicación Práctica}

En adquisiciones largas (por ejemplo, 25 horas), usar shift registers permite guardar datos progresivamente sin esperar al final del proceso.

---

\section{Programa 5 – Reconstrucción del String}

\subsection*{Objetivo}
Reforzar el uso de loops y concatenación.

\subsection*{Procedimiento}

\begin{enumerate}
    \item Crear un nuevo FOR Loop.
    \item Iterar sobre el arreglo de letras.
    \item Utilizar:
    \begin{itemize}
        \item \textbf{Concatenate Strings}
        \item Shift Register para acumulación
    \end{itemize}
\end{enumerate}

\subsection*{Proceso}

\begin{center}
\texttt{"" $\rightarrow$ "h" $\rightarrow$ "he" $\rightarrow$ "hel" $\rightarrow$ ...}
\end{center}

Hasta reconstruir:

\begin{center}
\texttt{hello world}
\end{center}

---

\section{Programa 6 – Extracción de Texto y Número}

\subsection*{Objetivo}
Aprender a convertir datos tipo string en valores numéricos.

\subsection*{Procedimiento}

\begin{enumerate}
    \item Modificar el string final a:
    \begin{center}
    \texttt{hello world 10}
    \end{center}
    \item Insertar el bloque \textbf{Scan From String}.
    \item Usar el formato:
    \begin{center}
    \texttt{\%s \%s \%d}
    \end{center}
\end{enumerate}

\subsection*{Resultado}

\begin{itemize}
    \item Salida 1: \texttt{hello}
    \item Salida 2: \texttt{world}
    \item Salida 3: \texttt{10} (convertido a entero)
\end{itemize}

El valor numérico puede utilizarse para operaciones matemáticas posteriores.

---

\section*{Conceptos Clave Aprendidos}

\begin{itemize}
    \item Manipulación básica de strings.
    \item Indexación y substrings.
    \item Uso de FOR Loops.
    \item Automatización según longitud.
    \item Construcción dinámica de arreglos.
    \item Diferencia entre shift register y auto-indexing.
    \item Parsing y conversión de tipos.
    \item Aplicación en adquisición de datos.
\end{itemize}

\end{document}